\subsection{Mối tương quan giữa các tuyến trong giao thông đồ thị}
Dựa trên những phân tích này, nhóm xây dựng mô hình Transformer encoder-only để học biểu
diễn không gian--thời gian cho từng vùng giao thông. Tuy nhiên, khác với các bài toán đồ thị tĩnh,
bài toán giao thông đô thị đòi hỏi mô hình phải nắm bắt được \textbf{tương tác động và có độ trễ theo
thời gian giữa các tuyến đường}, phản ánh cơ chế lan truyền của ùn tắc trong mạng lưới giao thông.

Để giải quyết yêu cầu này, nhóm sẽ khai thác các mối tương quan chéo có xét đến độ trễ thời gian
giữa các tuyến đường, dựa trên khuôn khổ mạng tương quan giao thông được đề xuất bởi Guo
et al.~\cite{guo2019identifying}. Cách tiếp cận này cho phép định lượng mức độ ảnh hưởng giữa
các tuyến đường và xác định các tuyến có vai trò quan trọng trong quá trình lan truyền ùn tắc.

Cụ thể, với mỗi cặp tuyến đường $i$ và $j$, từ chuỗi thời gian vận tốc đã được chuẩn hóa
$S_i(t)$ và $S_j(t)$, hàm tương quan chéo có độ trễ được định nghĩa như sau:
\begin{equation}
X_{i,j}(\tau) =
\frac{
\sum_{t=1}^{L-\tau} (S_i(t) - \bar{S}_i)\,(S_j(t+\tau) - \bar{S}_j)
}{
\sqrt{
\sum_{t=1}^{L-\tau} (S_i(t) - \bar{S}_i)^2
\sum_{t=1}^{L-\tau} (S_j(t+\tau) - \bar{S}_j)^2
}
},
\end{equation}
trong đó $\tau \in [-\tau_{\max}, \tau_{\max}]$ là độ trễ thời gian, $L$ là độ dài chuỗi thời gian,
và $\bar{S}_i$, $\bar{S}_j$ lần lượt là giá trị trung bình của các chuỗi tương ứng. Với $\tau < 0$,
ta có $X_{i,j}(\tau) = X_{j,i}(-\tau)$.

Từ hàm tương quan chéo này, Guo et al.~\cite{guo2019identifying} đề xuất trọng số liên kết dương
(positive link weight) giữa hai tuyến đường $i$ và $j$ được xác định bởi:
\begin{equation}
W^{\text{pos}}_{i,j} =
\frac{
\max_{\tau} X_{i,j}(\tau) - \text{mean}_{\tau} \left( X_{i,j}(\tau) \right)
}{
\text{std}_{\tau} \left( X_{i,j}(\tau) \right)
},
\end{equation}
trong đó $\max_{\tau}$, $\text{mean}_{\tau}$ và $\text{std}_{\tau}$ lần lượt là giá trị lớn nhất,
trung bình và độ lệch chuẩn của hàm tương quan chéo trên miền $\tau$. Đồng thời, độ trễ tương ứng
với ảnh hưởng mạnh nhất được xác định là:
\begin{equation}
\tau^{\text{pos}}_{i,j} = \arg\max_{\tau} X_{i,j}(\tau).
\end{equation}

Trọng số $W^{\text{pos}}_{i,j}$ phản ánh mức độ nổi trội của đỉnh tương quan so với nền nhiễu,
giúp lọc bỏ các tương quan yếu hoặc ngẫu nhiên, trong khi $\tau^{\text{pos}}_{i,j}$ cung cấp thông
tin về hướng và độ trễ của ảnh hưởng giữa các tuyến đường. Việc sử dụng các đại lượng này cho
phép nhóm xác định các liên kết giao thông mang ý nghĩa chức năng, phù hợp để mô hình hóa
quá trình lan truyền ùn tắc trong không gian và thời gian.

Trên cơ sở các liên kết tương quan có ý nghĩa thống kê và ràng buộc không gian, nhóm xây
dựng các vùng giao thông chức năng và xác định tập các tuyến mục tiêu cần dự đoán. Các vùng
này sau đó được đưa vào mô hình Transformer encoder-only với các cơ chế mã hóa không gian--
thời gian, cho phép học biểu diễn giàu ngữ cảnh và nâng cao khả năng dự báo ùn tắc giao thông
trong môi trường đô thị phức tạp.
